\RequirePackage{fontspec}
%\setmainfont{Latin Modern Roman}
\usepackage[brazilian]{babel} % 28/04/2026 Idioma principal
\usepackage[babel,final]{microtype} % Justificação (compilação final)
\usepackage[cmyk]{xcolor} % Cores cmyk (impressão)
\usepackage[default]{fontsetup}
\usepackage[none]{hyphenat} % Remover hifenizações
\usepackage[useregional]{datetime2} % Data no formato do idioma do documento

\usepackage{calc, parselines} % Símbolos matemáticos simples e processamento do texto 13/06/2026
\usepackage{csquotes, unicode-math}

\usepackage{etoolbox} % for \ifstrempty e final dos livros.
\usepackage{setspace} % Espaçamento entre linhas
\usepackage{fancybox, graphicx} % Caixas de texto e figuras/gráficos 03/05/2026
\usepackage{geometry, multicol} % dimensões, colunas
\usepackage{hologo, url} % TeX logos e urls.
\usepackage{balance} % colunas balanceadas
% \usepackage[titles]{tocloft}
%\usepackage[utf8]{inputenc}

% Cores
\definecolor{Black}{cmyk}{1, 1, 1, 1} % Preto puro
\definecolor{Cinza}{cmyk}{0, 0, 0, 0.85}
\definecolor{AzulMarinho}{cmyk}{1, 1, 0, 0.5}
\definecolor{Azul}{cmyk}{1, 1, 0, 0.3}
%\definecolor{VermelhoVivo}{cmyk}{0, 1, 1, 0.2} % Original
%\definecolor{VermelhoVivo}{cmyk}{0, 1, 1, 0.25} % Opção 1
\definecolor{VermelhoVivo}{cmyk}{0, 1, 1, 0.3} % Em uso
\definecolor{Verde}{cmyk}{1, 0, 1, 0.55} % Original
%\definecolor{Verde}{cmyk}{1, 0, 1, 0.3} % Opção 1
%\definecolor{Magenta}{cmyk}{0, 0.95, 0.49, 0.05} % Original
\definecolor{Magenta}{cmyk}{0, 1, 0.5, 0.05} % Original
%\definecolor{Roxo}{cmyk}{0, 0.9, 0, 0.37} % Original
\definecolor{Roxo}{cmyk}{0, 0.9, 0, 0.5} % Original
% \definecolor{Laranja}{cmyk}{0, 0.6, 1, 0} % Em uso
\definecolor{Laranja}{cmyk}{0, 0.65, 1, 0.2}
% \definecolor{Marrom}{cmyk}{0, 0.7, 0.9, 0.5} % Original
%\definecolor{Marrom}{cmyk}{0, 0.65, 0.9, 0.61} % Em uso
\definecolor{Marrom}{cmyk}{0, 0.65, 0.95, 0.7} % Opção 1
\definecolor{Verse}{cmyk}{0, 1, 0.65, 0.25} % Cor do versículo - Púrpura
%\definecolor{Verse}{cmyk}{0.5, 0.5, 0.5, 0.7} % Cor do versículo - Cinza
\definecolor{Amarelo}{cmyk}{0, 0.2, 1, 0.1}
\definecolor{Oliva}{cmyk}{0, 0, 1, 0.5}

\newcommand{\Black}{\color{Black}}
\newcommand{\TextColor}{\Black}
\newcommand{\IntroColor}{\color{AzulMarinho}}
\newcommand{\OliveColor}{\color{Oliva}}
\newcommand{\WJ}{\color{VermelhoVivo}}
\newcommand{\Red}{\WJ}
\newcommand{\Blue}{\IntroColor}
\newcommand{\LightBlue}{\color{Azul}}
\newcommand{\Green}{\color{Verde}}
\newcommand{\YellowColor}{\color{Amarelo}}
\newcommand{\VerseColor}{\color{Verse}}

% Cores da legenda
\newcommand{\blue}[1]{{\color{Azul}#1}}
\newcommand{\red}[1]{{\color{VermelhoVivo}#1}}
\newcommand{\green}[1]{{\color{Verde}#1}}
\newcommand{\magenta}[1]{{\color{Magenta}#1}}
\newcommand{\purple}[1]{{\color{Roxo}#1}}
\newcommand{\orange}[1]{{\color{Laranja}#1}}
\newcommand{\brown}[1]{{\color{Marrom}#1}}

% Título curto em coluna única e entrada no sumário, conteúdo em coluna dupla e
% última página balanceada
<<<<<<< HEAD
\newcommand{\ShortTitle}[1]{
%\flushcolsend
\clearpage\gdef\theshorttitle{#1}\gdef\webbook{#1}\twocolumn[\vspace{-6ex}\section*{}\addcontentsline{toc}{section}{#1}\begin{center}{\LARGE\centering\bfseries #1 \par}\end{center}]
%\flushend
} % Funcional
=======
\newcommand{\ShortTitle}[1]{\flushcolsend\clearpage\gdef\theshorttitle{#1}\gdef\webbook{#1}\twocolumn[\vspace{-6ex}\section*{}\addcontentsline{toc}{section}{#1}\begin{center}{\Large\centering\bfseries #1 \par}\end{center}]\flushend}
>>>>>>> origin/lualatex

\newcommand{\theshorttitle}{} % inicializa

\newcommand{\PeriphTitle}[1]{%
\clearpage
\gdef\theshorttitle{#1}%
\gdef\webbook{#1}%
\NormalFont
}

% Título principal
\newcommand{\MT}{\centering\Large\bfseries\Black} % Funcional

% Capítulo comum
\newcommand{\Chap}[1]{{\Large\bfseries\vspace{1.78ex}\centering\setcounterpageref{footnote}{0}\parindent=0mm #1\nopagebreak\par}\def\webchap{#1:}\nopagebreak}
\newcommand{\OneChap}{\vspace{1.78ex}\def\webchap{}}
\newcommand{\ChapOne}[1]{\vspace{1.78ex}\def\webchap{#1:}}
\newcommand{\PsalmChap}[1]{{\vspace{1.78ex}\centering #1\nopagebreak\par}\def\webchap{#1:}\nopagebreak}
\newcommand{\PoetryChap}[1]{{\vspace{1.78ex}\centering #1\nopagebreak\par}\def\webchap{#1:}\nopagebreak}

% Parágrafo de abertura de capítulo
\newcommand{\PP}{\par\noindent}

% Versículo inicial do capítulo
\newcommand{\VerseOne}[1]{\VS{#1}}
\newcommand{\VP}[1]{\VS{#1}}
\newcommand{\VA}[1]{\VS{#1}}
\newcommand{\FV}[1]{\VS{#1}}

% Versículo padrão
\newcommand{\VS}[1]{{\TextColor\raisebox{\RV}{{\footnotesize#1}}\hspace{1mm}\markboth{\webbook\ \webchap #1}{\webbook\ \webchap #1}}} % Verse number

% Opções Nativas da classe Memoir (substituir footmisc/footkmpj)
\raggedbottom % Espaçamento flexível
\feettobottom % Posição no rodapé
\feetbelowfloat % Notas abaixo
\setlength{\footmarkwidth}{0.8em} % Largura da caixa da marca da nota
\setlength{\footmarksep}{0.2em} % Separação da marca e do texto.
\setlength{\footparident}{0em} % Recuo do parágrafo.
\renewcommand{\foottextfont}{\fontsize{10pt}{10pt}\selectfont} % Fonte da nota 10pt.

\makeatletter

<<<<<<< HEAD
% Garantia incondicional: injeta \\vfill\\relax antes do rodapé em qualquer página.
\renewcommand\*{\m@mopfootnote}{\setbox\@outputbox \vbox{%
  \boxmaxdepth\@maxdepth
  \@tempdima\dp\@cclv 
  \unvbox\@cclv%
  \vskip-\@tempdima%
  \vfill\relax % Expande o espaço em branco antes da nota
  \vskip \skip\footins%
  \color@begingroup%
    \normalcolor%
    \footnoterule%
    \unvbox \footins%
  \color@endgroup}}
=======
% % Inicia o multicols no começo do sumário (executado uma vez na primeira compilação)
% \AtBeginDocument{%
%   \addtocontents{toc}{\string\begin{multicols}{3}}%
% }

% % Fecha o multicols atual ANTES de escrever a linha da parte
% \pretocmd{\@part}{%
%   \addtocontents{toc}{\string\end{multicols}}%
% }{}{}
%
% % Abre um novo multicols DEPOIS de escrever a linha da parte
% \apptocmd{\@part}{%
%   \addtocontents{toc}{\string\begin{multicols}{3}}%
% }{}{}
%
% % Fecha o último bloco multicols no final do documento
% \AtEndDocument{%
%   \addtocontents{toc}{\string\end{multicols}}%
% }
>>>>>>> origin/lualatex

% Nota de rodapé com símbolo personalizado
\renewcommand{\thefootnote}{\small\alph{footnote})} % letras

% Formatação da marca
\renewcommand{\@makefnmark}{%
%\hbox{\raisebox{0.5ex}{\textcolor{footcolor}{\@thefnmark)}}} % Anterior
\hbox{\raisebox{0.5ex}{\textcolor{footcolor}{\@thefnmark}}} % 
}

% Define FTNT, uma cor 'footcolor' com a cor atual etc.
\newcommand{\FTNT}[2]{\colorlet{footcolor}{.}\footnote{\textcolor{footcolor}{\hspace{0.5mm}#2}}}

\newcommand{\XREF}[2]{
\footnote{\textcolor{Black}{
\bfseries\itshape#2}}}

\makeatother

% Referência para rodapé (Negrito)
\newcommand{\FR}[1]{\bfseries#1}

% Adição de texto (ex.: {\\ADD{...}})
\newcommand{\ADD}[1]{\textit{#1}}

% Opcional: remover pontos dos líderes
%\renewcommand{\cftchapleader}{\cftdotfill{\cftdotsep}}
%\renewcommand{\cftsecleader}{\cftdotfill{\cftdotsep}}

% Fonte normal
\newcommand{\NormalFont}{\normalfont}
\NormalFont
\newcommand{\SmallFont}{\small}
